\section[Fargues--Fontaine 曲线 I]{Fargues--Fontaine 曲线 I（11 月 2 日）}
\srcpage{7}

\subsection{局部域}

固定一个非阿基米德局部域 $E$，其剩余域是特征为 $p$ 的 $\mathbb{F}_q$。设 $\pi\subset\mathcal{O}_E$ 为一个一致化子（uniformizer）\yizhu{原文如此，此处应为 $\pi\in\mathcal{O}_E$。}{A（原文 §2.1 p.7）}。

\begin{example}
具体地说，或者 $E\cong\mathbb{F}_q((t))$，或者 $E$ 是 $\mathbb{Q}_p$ 的有限扩张。
\end{example}

目标是「把 $\Spec E$ 几何化」。

\begin{itemize}
  \item $\Spec E$ 的 Zariski 网站是一个点，这不太几何——我们要把它加强。
  \item $E$ 的 étale 网站是 $\{\text{有限可分}\ E\text{-代数}\}^{\mathrm{op}}$。换言之，它是 $B\Gal(\overline{E}/E)$，
        也可以看作带连续 $\Gal(\overline{E}/E)$-作用的有限集合范畴。这更有趣，但仍然不太几何。
\end{itemize}

我们有一个短正合列
\[
\begin{tikzcd}[column sep=2.4em, row sep=1.4em]
0 \arrow[r] & I_E \arrow[r] & \Gal(\overline{E}/E) \arrow[r]
  & \Gal(\overline{\mathbb{F}}_q/\mathbb{F}_q) \arrow[r] \arrow[d, "\sim"] & 1 \\
  & & & \widehat{\mathbb{Z}}=\langle\Frob_q\rangle &
\end{tikzcd}
\]
并且惯性子群 $I_E$ 还可以进一步分解为
\[
0\to P_E\to I_E\to\prod_{\ell\ne p}\mathbb{Z}_\ell\to0 .
\]

\subsubsection{局部 Tate 对偶}

若 $E$ 的特征为 $0$，则对一切挠 $\Gal(\overline{E}/E)$-表示 $M$，有配对
\[
H^i_{\et}(\Spec E,M)\otimes H^{2-i}_{\et}(\Spec E,M^*(1))
\longrightarrow H^2_{\et}(\Spec E,\mathbb{Q}/\mathbb{Z}(1))=\mathbb{Q}/\mathbb{Z}.
\]
这里 $M^*=\Hom(M,\mathbb{Q}/\mathbb{Z})$，而 $\mathbb{Q}/\mathbb{Z}(1)$ 是 Tate 扭，即 $\bigcup_n\mu_n$。
若 $E$ 的特征为 $p$，则我们把讨论限制在 prime-to-$p$ 的 $M$ 与 $\mathbb{Q}/\mathbb{Z}$ 中的 prime-to-$p$ 挠部分上。

这看起来像 2 维流形上的 Poincaré 对偶。（但 Tate 扭意味着它的定向层非平凡。）
我们想把 $\Spec E$ 变成更接近紧 Riemann 曲面的东西。

\subsection{等特征的 Fargues--Fontaine 曲线}

设 $E=\mathbb{F}_q((t))$，并令 $\breve E=\overline{\mathbb{F}}_q((t))$。直观上，$\Spec\overline{\mathbb{F}}_q((t))$ 是
「$\overline{\mathbb{F}}_q$ 上的形式去心开单位圆盘（formal punctured open unit disc）」。但它当然只有一个点。
我们将通过过渡到扩域来「造出更多空间」。

取 $C/\overline{\mathbb{F}}_q$ 为完备代数闭非阿基米德域，例如 $C=\overline{\mathbb{F}}_q((u))^\wedge$。考虑 adic 谱
$\Spa C\times_{\Spa\mathbb{F}_q}\Spa\mathbb{F}_q((t))$。这就是「$C$ 上的去心开单位圆盘」
\[
\mathbb{D}^*_C=\{x\mid 0<|x|<1\}.
\]
它现在有许多点。

基变换到 $C$ 相当于过渡到「几何」的情形；正如前几讲那样，我们要记住 Frobenius $\varphi_C$。

\begin{definition}
（关于 $E,C$ 的）\emph{Fargues--Fontaine 曲线}\footnote{这依赖于 $C$ 的选取。}是
$X_{C,E}:=\mathbb{D}^*_C/\varphi_C^{\mathbb{Z}}$，视作 $E$ 上的 adic 空间。
\end{definition}

\srcpage{8}

\begin{remark}
$\varphi_C^{\mathbb{Z}}$ 在 $\mathbb{D}^*_C$ 上的作用自由。因此这个商的性质非常好。
然而要注意：尽管 $\mathbb{D}^*_C$ 是 $C$ 上的有限型对象，由于作用在底 $C$ 上非平凡，它不再生活在 $C$ 上；
它反而是 $E$ 上的无限型对象。
\end{remark}

这在何种意义上是对 $\Spec E$ 的几何化？下面这些性质归功于 Fargues--Fontaine。
\begin{itemize}
  \item $H^0(X_{C,E},\mathcal{O}_{X_{C,E}})=E$。
  \item $X_{C,E}$ 的有限 étale 覆盖的范畴是 $(\Spec E)_{\et}$。
  \item $H^i_{\et}(X_{C,E},M)\cong H^i_{\et}(\Spec E,M)$。
\end{itemize}
特别地，局部 Tate 对偶变成 $X_{C,E}$ 上的整体对偶。

\subsection{关于 adic 空间的一些回顾}

粗略地说，adic 空间是联系于某些拓扑环（例如 Banach 代数）的概形变体。它们与概形的共同之处在于：
\begin{itemize}
  \item 点之间有特殊化关系。
  \item 它们非常一般：基础部分没有任何有限性假设（这点与刚性解析几何不同）。
\end{itemize}

\begin{definition}
设 $A$ 是拓扑环。
\begin{enumerate}[label=(\arabic*), leftmargin=2.6em]
  \item 称 $A$ 是 \emph{adic} 的，若存在理想 $I\subseteq A$，使得 $\{I^n\mid n\ge0\}$ 是 $0$ 的邻域基。
        这样的理想 $I$ 称为\emph{定义理想}（ideal of definition）。
        （它并不唯一；但对任意两个 $I,J$，存在 $n$ 使得 $I^n\subset J$ 且 $J^n\subset I$。）
  \item 称 $A$ 是 \emph{Huber} 的\footnote{在 Huber 的原始论文中，这被称为「$f$-adic」。}，
        若存在开子环 $A_0\subset A$，它是定义理想为有限生成的 adic 环。
        这样的 $A_0\subseteq A$ 称为\emph{定义环}（ring of definition）。
\end{enumerate}
\end{definition}

\begin{remark}
任何这样的 $A$ 都有完备化 $\widehat{A}$，它包含 $I$-adic 完备化 $\widehat{A}_0\subset\widehat{A}$ 作为开子环。
\end{remark}

对我们最重要的一类 adic 环是：

\begin{definition}
若 $A$ 含有一个拓扑幂零单位 $\varpi\in A$，则称 $A$ 是 \emph{Tate 环}。
这样的 $\varpi$ 称为\emph{伪一致化子}（pseudouniformizer）。
\end{definition}

\begin{example}
任何非阿基米德域都是 Tate 的。在 $\mathbb{F}_p((t))$ 中可取 $\varpi=t$；在 $\mathbb{Q}_p$ 中可取 $\varpi=p$。
更一般地，非阿基米德局部域上的任何 Huber 环都是 Tate 的。
\end{example}

\begin{remark}
设 $K$ 是任意非阿基米德局部域，$\varpi\in K$ 是伪一致化子，$A/K$ 是完备 Huber 环。
则 $A$ 自然地具有 $K$ 上 Banach 代数的结构，其「单位球」为
$\{f\in A:|f_0|\le1\}=A_0$\yizhu{原文下标作 $|f_0|$，疑为笔误，当作 $|f|\le1$。}{A（原文 §2.2 p.8）}。
任何定义环 $A_0\subset A$ 都带有 $\varpi$-adic 拓扑。
\end{remark}

$A$ 上的范数可以这样构造：任取 $\|\varpi\|$ 为 $(0,1)$ 中的某个实数，然后定义
\[
\|a\|=\inf_{\{n\mid\varpi^na\in A_0\}}2^n .
\]

$K$ 上带连续映射的 Banach 代数与 $K$ 上的 Tate--Huber 环之间有范畴等价。
所以我们本可以用 Banach 代数而不是 Tate 环来表述。但 Tate 环的进路至少有两点好处：
\begin{itemize}
  \item 我们不需要固定一个非阿基米德的基域 $K$。
  \item \srcpage{9}范数本身是多余的结构。
\end{itemize}

\begin{definition}
Huber 环 $A$ 的\emph{赋值谱}（valuation spectrum）是
\[
\mathrm{Cont}(A):=\{\text{连续赋值}\ |\cdot|:A\to\Gamma\cup\{0\}\}/\sim,
\]
其中 $\Gamma$ 是全序交换群（例如 $\mathbb{R}_{>0}$）。
我们在 $\mathrm{Cont}(A)$ 上取由开集 $\{|f|\le|g|\ne0\}\subset\mathrm{Cont}(A)$（$f,g\in A$）生成的拓扑。
\end{definition}

具体来说，「连续赋值」是指：
\begin{itemize}
  \item $|ab|=|a|\cdot|b|$，$0\gamma=\gamma0=0$，且对一切 $\gamma$ 有 $0\le\gamma$。
  \item $|a+b|\le\max(|a|,|b|)$。
  \item $|0|=0$，$|1|=1$。
  \item 对一切 $\gamma\in\Gamma$，$\{a\mid |a|<\gamma\}\subset A$ 是开集。
\end{itemize}

两个连续赋值 $|\cdot|_1,|\cdot|_2$ 称为\emph{等价}的，若它们给出的二元关系相同：
\[
|a|_1\ge|b|_1\iff|a|_2\ge|b|_2 .
\]
这等价于：若 $\Gamma_i$ 取得极小，则存在同构 $\Gamma_1\xrightarrow{\sim}\Gamma_2$，使得下图交换。
\[
\begin{tikzcd}[column sep=3.2em, row sep=1.8em]
  & \Gamma_1\cup\{0\} \arrow[dd, "\sim"] & \\
A \arrow[ru, "|\cdot|_1"] \arrow[rd, "|\cdot|_2"'] & & \\
  & \Gamma_2\cup\{0\} &
\end{tikzcd}
\]

\begin{example}
若 $\Gamma$ 可嵌入 $\mathbb{R}_{>0}$，则称该赋值是「秩 $1$」的。
一个自然出现、且不是秩 $1$ 的 $\Gamma$ 的例子是 $\mathbb{R}_{>0}\times\gamma^{\mathbb{Z}}$，
其中对一切 $r\in\mathbb{R}_{>1}$ 都有 $r>\gamma>1$。
\end{example}

\begin{example}
定义中的严格性只在 $\Gamma=\{1\}$ 时才起作用。
\end{example}

\begin{definition}
一个 \emph{Huber 对}（Huber pair）是一对 $(A,A^+)$，其中 $A$ 是 Huber 环，
$A^+\subset A$ 是由幂有界（power-bounded）元素构成的、开的整闭子环。
记 $A^\circ\subset A$ 为幂有界元素构成的子环。
\end{definition}

我们定义
\[
\Spa(A,A^+)=\{|\cdot|:|A^+|\le1\}\subset\mathrm{Cont}(A),\qquad
\Spa A=\Spa(A,A^\circ).
\]

我们可以在有理子集组成的基上，给 $\Spa A$ 配备 Huber 环的预层
$\mathcal{O}_{\Spa A}\supset\mathcal{O}^+_{\Spa A}$。这里的有理子集（rational subsets）形如
\[
U\Bigl(\frac{f_1,\dots,f_n}{g}\Bigr):=\bigl\{|f_i|\le|g|\ne0\bigr\},
\]
其中 $(f_1,\dots,f_n,g)$ 生成一个开理想。于是
\[
\mathcal{O}_{\Spa A}\Bigl(U\Bigl(\frac{f_1,\dots,f_n}{g}\Bigr)\Bigr)
=A\Bigl\langle\frac{f_1}{g},\dots,\frac{f_n}{g}\Bigr\rangle
\]
「允许 $f_i/g$ 的一切收敛级数」。$+$ 子环的作用是记录哪些元素 $\le1$，所以它包含 $f_i/g$。
子环 $\mathcal{O}^+_{\Spa A}\bigl(U(\frac{f_1,\dots,f_n}{g})\bigr)$ 是包含 $A^+$ 与所有 $f_i/g$ 的最小子环。

\begin{theorem}[Huber 等]
在一切实际情形中，$\mathcal{O}_{\Spa(A,A^+)}$ 都是层，但并非总是如此。
\end{theorem}

\srcpage{10}
\begin{center}
\figP{010}
\end{center}

\begin{remark}
最近人们认识到（Bambozzi--Kremnizer、Clausen--S），这种非层性可以通过允许结构层取导出对象来修正。
微妙之处在于如何把导出色性与拓扑结合起来；我们与 Clausen 借助凝聚态（condensed）数学的理论解决了这一点。
这与本课程无关。
\end{remark}

\begin{definition}
一个 \emph{adic 空间}是一个三元组 $(X,\mathcal{O}_X,\mathcal{O}^+_X)$，其中
\begin{itemize}
  \item $X$ 是拓扑空间，
  \item $\mathcal{O}_X$ 是完备拓扑环的层，
  \item $\mathcal{O}^+_X$ 是 $\mathcal{O}_X$ 的子层，
\end{itemize}
它在局部上形如 $(\Spa(A,A^+),\mathcal{O}_A,\mathcal{O}^+_A)$，再加上「局部环化」的正确类比。
\end{definition}

\begin{remark}
这也可以等价地表述为：给出 $\mathcal{O}_X$ 的茎上的赋值，而不是给出 $\mathcal{O}^+_X$。
子层 $\mathcal{O}^+_X\subset\mathcal{O}_X$ 可以还原为那些在所有点处的赋值都 $\le1$ 的元素。
\end{remark}

\begin{remark}
为什么我们要求 $\mathcal{O}_X$ 是完备环的层？一个好处是：只有借助完备的对象，
我们才能表述有理子集上函数 $A\langle\frac{f_1}{g},\dots,\frac{f_n}{g}\rangle$ 的良好泛性质。
\end{remark}

\begin{example}
设 $C$ 是非阿基米德域。我们来解释 $\mathbb{D}^*_C=\{x\mid0<|x|<1\}$ 上的 adic 空间结构。有 Tate 代数
\[
C\langle T\rangle=\Bigl\{\sum a_nT^n \mid a_n\in C,\ a_n\to0\Bigr\}.
\]
于是 $\Spa C\langle T\rangle=:\mathbb{B}_C$ 是 $C$ 上的「闭单位圆盘」$\{x\mid0\le|x|\le1\}$。
我们有 $\mathbb{B}^*_C=\mathbb{B}_C\setminus\{0\}=\bigcup_{\epsilon>0}\mathbb{A}(\epsilon,1)$，
其中非正式地说，对 $\epsilon\in|C|$，$\mathbb{A}(\epsilon,1)$ 就是环带 $\{x\mid\epsilon\le|x|\le1\}$。
更形式地说，$\mathbb{A}(\epsilon,1)=\Spa A_\epsilon$，其中
\[
A_\epsilon=\Bigl\{\sum_{n=-\infty}^{\infty}a_nT^n:\ a_n\in C,\
|a_n|\xrightarrow[n\to\infty]{}0,\ \epsilon^n|a_n|\xrightarrow[n\to-\infty]{}0\Bigr\}.
\]
类似地，$\mathbb{D}^*_C=\bigcup_{\epsilon>0,\,r<1}\mathbb{A}(\epsilon,r)$，其中
$\mathbb{A}(\epsilon,r)=\{\epsilon\le|T|\le r\}$。
\end{example}

\begin{warning}
$\mathbb{D}^*_C\subset\mathbb{B}_C$ 是开的，但它不等于 $\{0<|T|<1\}$。
存在一个点 $x\in\mathbb{B}_C=\Spa\langle T\rangle$，使得对一切 $r\in|C|$（$r<1$）都有 $r<|T(x)|<1$，
即 $x$ 无穷小地小于 $1$。递增的环带并集不包含这样的点。
子集 $\{0<|T|<1\}$ 不是开的，因为严格不等式 $<1$ 定义的是闭子集。
\end{warning}

\srcpage{11}
\begin{center}
\figP{011-1}
\end{center}
\begin{center}
\figP{011-2}
\end{center}

事实上存在自然的映射 $\mathrm{rad}:\mathbb{B}_C\to[0,1]$ 与 $\mathrm{rad}:\mathbb{D}^*_C\to(0,1)$，
它们都把 $x$ 送到 $|T(\tilde x)|$。Frobenius $\varphi_C$ 把 $\mathbb{D}^*_C$ 从原点向外移动，即
$\mathrm{rad}\circ\varphi_C=\mathrm{rad}^{1/q}$。

这说明 $\varphi_C$ 在 $\mathbb{D}^*_C$ 上的作用自由且真不连续（properly discontinuous）。
因此可以定义 $X_{C,E}=\mathbb{D}^*_C/\varphi_C^{\mathbb{Z}}$。
这可以看作 $\mathbb{A}(r,r^{1/q})$ 模去经由 $\phi:\mathbb{A}(r,r)\cong\mathbb{A}(r^{1/q},r^{1/q})$
对边界环带所作的等同。

直觉上，这个商看起来有点像复环面。
